% ---------------------------------------------------------
\chapter{Implementierung}\label{chap:implementation}
% ---------------------------------------------------------


Abbildung \ref{umlcd:figure} zeigt ein Klassendiagramm welches mit dem
\LaTeX-Paket
\texttt{pgf-umlcd}\footnote{\url{https://ctan.org/pkg/pgf-umlcd},
  Letzter Zugriff: 19.08.2022} erstellt wurde und Abbildung
\ref{umlsd:example} zeigt ein Sequenzdiagramm welches mit dem
\LaTeX-Paket
\texttt{pgf-umlsd}\footnote{\url{https://ctan.org/pkg/pgf-umlsd},
  Letzter Zugriff: 19.08.2022} generiert wurde.


\begin{figure}[h]
  \centering
\begin{tikzpicture}
  \begin{abstractclass}[text width=5.5cm]{Figure}{3,0}
    \attribute{-name : String}
    \attribute{-hitPoints : int}
    \attribute{-strength : int}
    \attribute{-dexterity : int}

    \operation{+attack(Figure enemy): int}
    \operation{+\textit{causeDamage(): int}}
    \operation{+isDead(): boolean}
  \end{abstractclass}
  
  \begin{class}[text width=3.5cm]{Hero}{0,-7}
    \attribute{-w: Weapon}
    \inherit{Figure}
  \end{class}

   
  \begin{class}[text width=3.5cm]{Monster}{6,-7}
    \inherit{Figure}
  \end{class}
\end{tikzpicture}
\caption{UML-Klassendiagramm für Spielfiguren.}
\label{umlcd:figure}
\end{figure}

\begin{figure}[h]
  \centering
\begin{sequencediagram}
\newthread{t}{:Thread}
\newinst[2]{foo}{:Foo}
\newinst[2]{bar}{:Bar}
\begin{call}{t}{gnarf()}{foo}{return}
\begin{call}{foo}{fnord()}{bar}{ return }
\end{call}
\end{call}
\begin{messcall}{t}{doAsynStuff()}{bar}
\end{messcall}
\end{sequencediagram}
\caption{Beispiel für ein Sequenzdiagramm.}
\label{umlsd:example}
\end{figure}



Der Client sendet Ping-Anfgragen, die vom Server mit einer
Pong-Nachricht beantwortet werden (siehe Abbildung
\ref{umlsd:clientrequest}).


\begin{figure}[h]
    \centering
  \begin{sequencediagram}
    \newinst[1]{c}{Client}
    \newinst[3]{s}{Server}
    \mess[1]{c}{ping}{s}
    \mess[1]{s}{pong}{c}
  \end{sequencediagram}
  \caption{Clientanfrage}
  \label{umlsd:clientrequest}
\end{figure}



